% Generated draft: requires mathematical and source review.
\documentclass[11pt]{article}
\usepackage[margin=25mm]{geometry}
\usepackage{fontspec}
\setmainfont{Latin Modern Roman}
\usepackage{amsmath,amssymb,mathtools,unicode-math}
\setmathfont{Latin Modern Math}
\usepackage{enumitem,needspace,xurl,hyperref}
\setlength{\parindent}{0pt}
\setlength{\parskip}{5pt}
\setlength{\emergencystretch}{3em}
\newcommand{\N}{\mathbb{N}}
\newcommand{\Z}{\mathbb{Z}}
\newcommand{\Q}{\mathbb{Q}}
\newcommand{\R}{\mathbb{R}}
\newcommand{\C}{\mathbb{C}}
\newcommand{\F}{\mathbb{F}}
\newcommand{\A}{\mathbb{A}}
\newcommand{\PP}{\mathbb{P}}
\newcommand{\E}{\mathbb{E}}
\newcommand{\eps}{\varepsilon}
\newcommand{\dd}{\,\mathrm{d}}
\begin{document}
\section*{UnsolvedMath ID: 11000070}
\textbf{Source problem code:} catalogue entry
\subsection{Definitions and mathematical statement}
Let $X_n$ be an $n$-dimensional Euclidean building, equipped with its usual CAT$(0)$ metric, and let $\psi:\operatorname{Mod}_g\to\operatorname{Isom}(X_n)$ be an isometric action. Problem 6.1 asks to determine all such actions when $n$ is sufficiently small compared with $g$. In particular, it asks for conditions guaranteeing
\[
 \exists x\in X_n\quad\forall h\in\operatorname{Mod}_g,
 \qquad \psi(h)x=x.
\]
Neither the dimension/genus bound nor the equivalence relation for classifying actions is prescribed. The genus and surface conventions, and possible completeness, local finiteness, thickness, type-preservation or semisimplicity assumptions on the building/action, are absent.

The source proposes a Helly technique attributed to [Fa4]. For a Dehn twist $T_\alpha$ it writes $F_\alpha=\operatorname{Fix}(\psi(T_\alpha))$ and reports these sets nonempty and convex. Fixed sets of CAT$(0)$ isometries are convex, but their nonemptiness in this setting requires the applicable hypotheses. The nerve has a simplex for each finite family with nonempty common intersection. The curve complex $\mathcal C_g$ has vertices represented by essential simple closed curves and, in the usual nonexceptional convention, simplices represented by disjoint families.

The source describes a map $\mathcal C_g\to X_n$, invokes the wedge-of-spheres homotopy type of the curve complex ([Iv1]), and suggests that filling these spheres in the contractible building should yield additional common fixed points. This is an outlined strategy: a map to a contractible target alone does not imply such intersections. The nerve construction, the dimension/Helly hypotheses, and the explicit sphere representatives needed for the argument remain to be supplied. The source specifically asks for a more explicit understanding of those spheres as part of this approach.
\subsection{Short English statement}
Classify mapping-class-group actions on Euclidean buildings of small dimension relative to genus, especially by finding global-fixed-point criteria. The proposed nerve/Helly approach has unstated hypotheses and construction steps.
\subsection{Sources}
\begin{itemize}
\item Local registry: \url{data/unsolvedmath/problems.json}, permanent UnsolvedMath ID $11000070$, exact source code \texttt{catalogue entry}.
\item Farb, \emph{Problems on Mapping Class Groups and Related Topics}, Problem 6.1, PDF page 53: \url{https://www.math.uchicago.edu/~farb/papers/mcgbook.pdf}. Supplied access date: 2026-07-29; extraction: pdf-font-structure; page-disambiguated numbering among 329 headings.
\item Supplied reference labels: [Fa4] for the Helly technique and [Iv1] for the curve-complex homotopy type. Full bibliography entries are absent. The registry mentions partial rigidity results without exact references.
\end{itemize}
\end{document}
